% =============================================================================
% Chapter 3: Literature Survey
% NexaLearn Graduation Project — Cairo University, Faculty of Engineering
% =============================================================================

\chapter{Literature Survey}
\label{ch:literature-survey}

This chapter surveys the foundational technologies, architectural patterns, and academic literature underpinning the NexaLearn platform. The survey is organised into six thematic sections that correspond to the principal technical pillars of the system. Section~\ref{sec:ddd-clean-arch} reviews Domain-Driven Design (DDD) and Clean Architecture, establishing the decomposition strategy used to partition NexaLearn into fifteen bounded contexts with explicit dependency rules. Section~\ref{sec:event-driven-outbox} examines event-driven architecture and the transactional outbox pattern, which together guarantee reliable asynchronous integration across contexts and with external providers. Section~\ref{sec:ai-elearning} surveys recent advances in AI-assisted e-learning---specifically OpenAI Whisper for speech-to-text transcription and large language model (LLM) pipelines for automated quiz generation---detailing how these technologies reduce instructor workload while maintaining content quality through human-in-the-loop review. Section~\ref{sec:video-streaming} covers HTTP Live Streaming (HLS), the Mux video platform, and signed URL access control, which collectively enable adaptive-bitrate playback and fine-grained entitlement enforcement. Section~\ref{sec:comparative-study} presents a structured comparison of prior academic and industrial e-learning backend approaches, highlighting gaps in domain modelling, event reliability, and AI integration that NexaLearn addresses. Finally, Section~\ref{sec:implemented-approach} justifies the specific technology and pattern choices---NestJS, Prisma, the outbox pattern, and modular monolith deployment---synthesising the surveyed literature into a coherent architectural rationale.

% =============================================================================
% Section 3.1: Background on Domain-Driven Design and Clean Architecture
% =============================================================================

\section{Background on Domain-Driven Design and Clean Architecture}
\label{sec:ddd-clean-arch}

\subsection{Domain-Driven Design Fundamentals}

Domain-Driven Design (DDD), as originally formulated by Evans~\cite{evans2003ddd}, provides a set of strategic and tactical modelling patterns for managing complexity in software systems whose value derives from intricate domain logic. The central premise is that the \textit{domain model}---a software embodiment of the business concepts, rules, and invariants---must be the primary artefact driving design decisions, rather than implementation technology or database schemas. Evans introduces the concept of a \textbf{ubiquitous language}: a shared vocabulary used by domain experts, developers, and all project stakeholders, embedded consistently in code, diagrams, and documentation. This language prevents the translation errors that arise when business concepts are renamed or restructured between specification documents and implementation.

At the strategic level, DDD prescribes decomposing large problem spaces into \textbf{bounded contexts}---explicit boundaries within which a particular ubiquitous language applies. Each bounded context owns its domain model, database schema, and internal logic; communication with other contexts occurs through well-defined integration contracts. A \textbf{context map} documents the relationships between bounded contexts, specifying whether the relationship is upstream/downstream (customer-supplier), conformist (the downstream adopts the upstream model without change), or shared-kernel (two contexts share a subset of the domain model). This decomposition prevents the architectural erosion described in Section~\ref{subsubsec:bounded-context-erosion} by enforcing that code in one context cannot directly access repositories or entities owned by another.

At the tactical level, DDD defines a vocabulary of building blocks:

\begin{description}
  \item[\textbf{Entity.}] An object defined by its identity rather than its attributes. Two entities with identical field values are distinguishable by their identifier. In NexaLearn, \texttt{CourseEntity} and \texttt{QuizEntity} are entities whose identity persists across state changes.
  \item[\textbf{Value Object.}] An immutable object defined solely by the values of its attributes. \texttt{QuestionOption}, \texttt{Address}, or monetary amounts are typical value objects. Equality is structural rather than identity-based.
  \item[\textbf{Aggregate.}] A cluster of entities and value objects treated as a single consistency boundary. Aggregates have a root entity (the \textbf{aggregate root}) that enforces invariants and mediates all external access to the aggregate's internals. In NexaLearn, \texttt{QuizEntity} is the aggregate root for a graph that includes \texttt{QuestionEntity} children and \texttt{QuestionOption} value objects; all operations that add, remove, or reorder questions must go through the quiz aggregate root, ensuring invariants such as ``a quiz must have at least one question'' and ``question ordering within a quiz is contiguous.''
  \item[\textbf{Domain Event.}] A record of something that happened in the domain that domain experts care about. \texttt{EnrollmentActivated}, \texttt{LessonCompleted}, and \texttt{PaymentSettled} are domain events in NexaLearn. Events are immutable facts published to interested contexts or external subscribers.
  \item[\textbf{Repository.}] A mechanism for encapsulating storage, retrieval, and search behaviour, presenting a collection-like interface to the domain layer. Repositories per aggregate root abstract the underlying database technology---PostgreSQL, Redis, or in-memory---from domain logic.
  \item[\textbf{Domain Service.}] A stateless operation that does not naturally belong to an entity or value object. Domain services orchestrate multiple aggregates or coordinate with external systems while keeping business rules in the domain layer.
\end{description}

NexaLearn's codebase applies these tactical patterns systematically. The \texttt{QuizEntity} aggregate root, for example, owns a list of \texttt{questions}, enforces the invariant that minimum passing score is between 0 and 100, and exposes behaviour methods such as \texttt{addQuestion()} rather than allowing direct mutation of the questions array.

\subsection{Clean Architecture Layers}

Martin's Clean Architecture~\cite{martin2017clean} refines the layering concerns earlier expressed in the Hexagonal Architecture and Onion Architecture patterns. The defining principle is the \textbf{Dependency Rule}: source code dependencies must point \textit{inward}, from outer layers toward inner layers. Nothing in an inner layer can know about anything in an outer layer.

Applied to NexaLearn, the four layers are:

\begin{enumerate}
  \item \textbf{Domain Layer (innermost).} Contains entities, value objects, aggregate roots, domain events, and repository \textit{interfaces} (ports). This layer has zero dependencies on frameworks, databases, or external libraries. The ubiquitous language lives here: \texttt{Course}, \texttt{Enrollment}, \texttt{Payment}, \texttt{VideoAsset} are expressed as pure TypeScript classes with behaviour and invariants.
  \item \textbf{Application Layer.} Contains use-case interactors that orchestrate domain objects to fulfil a single user story. Each use case (e.g., \texttt{SubmitQuizAttemptUseCase}, \texttt{ActivateEnrollmentUseCase}) depends on repository interfaces from the domain layer but has no knowledge of how those repositories are implemented. The application layer may also define application-specific DTOs and events.
  \item \textbf{Infrastructure Layer.} Implements the repository interfaces declared in the domain layer. Prisma-based repository implementations, Mux API clients, Stripe SDK wrappers, and outbox event publishers all reside here. This is also where the transactional outbox publisher, job queues, and scheduled task processors live.
  \item \textbf{Presentation Layer (outermost).} Controllers, HTTP request/response DTOs, validation pipes, guards, and interceptors. NestJS controllers in NexaLearn are thin: they parse HTTP input, delegate to an application use case, and format the response. No business logic resides in controllers.
\end{enumerate}

The Dependency Rule is enforced through TypeScript's module system and NestJS module boundaries. A custom ESLint rule in the NexaLearn codebase disallows infrastructure-layer modules from importing any NestJS decorator or Prisma client code into domain or application layer files.

\subsection{Application to NexaLearn's Bounded Contexts}

NexaLearn decomposes its domain into fifteen bounded contexts, each mapped to a NestJS module (or set of related modules). The decomposition follows Dossena et al.'s~\cite{dossena2022ddd} recommendations for applying DDD to e-learning platforms, where enrolment, assessment, content management, and analytics each form natural bounded contexts because they involve distinct stakeholders, disparate lifecycle rules, and different persistence trade-offs. The \textbf{Courses} context, for example, owns course catalogue data, curriculum structure, and lesson content. It exposes read-only projections for the \textbf{Progress} context via materialised views and reacts to \texttt{EnrollmentActivated} events published by the \textbf{Enrollment} context. The \textbf{Assessment} context owns quizzes, questions, attempts, and scoring logic; its aggregate root (\texttt{QuizEntity}) enforces invariants such as maximum attempts per learner and question ordering. This bounded-context-per-module approach, combined with Clean Architecture layering within each module, yields a codebase where a developer can understand the Quiz domain entirely within the Assessment module without traversing cross-cutting infrastructure concerns.

The modular monolith deployment strategy (detailed in Section~\ref{sec:implemented-approach}) packages all fifteen contexts into a single NestJS application process, but the DDD boundaries and Clean Architecture layers ensure that the system remains maintainable as it grows. Each context can be extracted into an independent microservice with minimal refactoring if scaling or team autonomy requirements change, following the split-pattern described by Yang et al.~\cite{yang2023micro}.

Ultimately, DDD and Clean Architecture provide NexaLearn with three benefits directly motivated by the problem statement of Chapter~1. First, \textbf{testability}: domain entities with no framework dependencies can be unit-tested without NestJS test containers or database connections. Second, \textbf{maintainability}: a developer can reason about the Enrollment aggregate without understanding Stripe webhook processing details. Third, \textbf{bounded context isolation}: the Courses context can evolve its schema and domain logic without cascading changes to Discussion or Streak, provided the published event contracts remain stable.

% =============================================================================
% Section 3.2: Background on Event-Driven Architecture and the Outbox Pattern
% =============================================================================

\section{Background on Event-Driven Architecture and the Outbox Pattern}
\label{sec:event-driven-outbox}

\subsection{Event-Driven Architecture in Distributed Systems}

Event-Driven Architecture (EDA) is a software paradigm in which components communicate by producing and consuming \textit{events}---immutable records of something that has occurred, rather than by issuing direct command invocations~\cite{fowler2005eda}. In e-learning systems, events arise naturally from learner interactions (lesson viewed, quiz submitted, discussion post created), system state changes (enrolment activated, certificate issued), and external provider callbacks (payment succeeded, video ready). EDA decouples event producers from consumers: a single enrolment activation event may be consumed by the Progress context (to initialise progress projections), the Media context (to issue playback entitlements), the Notifications context (to send a welcome email), and the Search context (to update enrolment counts), all without the Enrollment context having any awareness of these downstream dependents.

The principal challenge in EDA is reliable event delivery. When a domain operation both modifies the database and publishes an event, these two actions must be atomic: either both succeed or neither takes effect. Performing the database write first, then publishing the event, creates a window in which the database write is visible but the event is not published---if the application crashes between the two operations, or if the message broker is unreachable, the system enters an inconsistent state. Conversely, publishing the event first and then writing the database risks acting upon an event whose associated state does not yet exist. This is known as the \textbf{dual-write problem}.

\subsection{The Transactional Outbox Pattern}

The transactional outbox pattern, popularised by Richardson~\cite{richardson2018outbox} in the microservices community and independently described as the ``application events'' pattern in messaging literature, solves the dual-write problem by persisting events in the same database transaction as the domain mutation. Rather than publishing an event immediately to a message broker, the application inserts a row into an \texttt{outbox\_events} table within the same database transaction that updates the domain aggregates. A separate background process---the \textbf{outbox publisher}---polls the outbox table for unprocessed events, publishes them to the message broker (or a message queue), and marks them as processed.

This design guarantees \textbf{at-least-once delivery}: an event may be published more than once if the publisher crashes after publishing but before marking the event as processed, but it will never be lost. Consumer applications use idempotency keys to achieve exactly-once semantics despite at-least-once transport.

\subsection{NexaLearn's Outbox Implementation}

NexaLearn implements a variant of the transactional outbox pattern using PostgreSQL as the event store and a dedicated NestJS \texttt{OutboxProcessorService} as the publisher. The key components are:

\begin{itemize}
  \item \textbf{OutboxEvent Entity.} A database table with columns: \texttt{id} (UUID primary key), \texttt{aggregate\_type}, \texttt{aggregate\_id}, \texttt{event\_type}, \texttt{payload} (JSONB), \texttt{created\_at}, \texttt{processed\_at} (nullable), and \texttt{retry\_count}.
  \item \textbf{OutboxPublisher Service.} A domain service that receives a domain event and inserts a corresponding outbox row. This service is invoked by domain services or use-case interactors after successfully modifying domain aggregates.
  \item \textbf{OutboxProcessorService.} A polling background worker (implemented as a NestJS \texttt{@Cron} scheduled task or a dedicated \texttt{@Processor} job queue consumer) that periodically queries \texttt{SELECT * FROM outbox\_events WHERE processed\_at IS NULL ORDER BY created\_at ASC LIMIT $n$ FOR UPDATE SKIP LOCKED}. The \texttt{FOR UPDATE SKIP LOCKED} clause ensures that multiple concurrent processor instances can work on different batches without contention---a critical property for horizontal scaling.
\end{itemize}

\subsection{Idempotency and Exactly-Once Processing}

Idempotency ensures that processing the same event multiple times produces the same outcome. NexaLearn's \textbf{ProjectionEventLedger} table records which events have been applied to each read model. The ledger schema includes: \texttt{id} (UUID), \texttt{event\_id} (foreign key to outbox\_events), \texttt{projection\_type} (e.g., \texttt{course-progress-view}), \texttt{handler\_id} (the specific handler instance that processed the event), and \texttt{applied\_at}. Before applying an event to a projection, the handler attempts \texttt{INSERT INTO projection\_event\_ledger (event\_id, projection\_type, handler\_id) VALUES (\$1, \$2, \$3) ON CONFLICT (event\_id, projection\_type) DO NOTHING}. If the insert succeeds, the event has not been processed and the handler proceeds; if it returns zero rows (due to the conflict), the event was already processed and the handler skips it.

This approach is inspired by the idempotency-key pattern used by Stripe~\cite{stripe2024api} for payment processing, adapted for internal event consumption. The combination of transactional outbox and idempotent projection handlers eliminates the dual-write problem without requiring distributed transactions, a two-phase commit coordinator, or an external message broker---a deliberate simplification discussed in Section~\ref{sec:implemented-approach}.

\subsection{CQRS-Lite Read Models}

While NexaLearn does not implement full Command Query Responsibility Segregation (CQRS)~\cite{fowler2011cqrs} with separate write and read databases, it adopts a ``CQRS-lite'' pattern: several bounded contexts maintain specialised read models that are updated asynchronously by outbox event handlers. These read models are optimised for specific query patterns that would require expensive joins or recomputation across the write-optimised domain schema. Examples include:

\begin{itemize}
  \item \textbf{CourseProgressView.} Denormalised per-learner progress summary showing total lessons, completed lessons, quiz scores, and streak data, updated reactively when \texttt{LessonCompleted}, \texttt{QuizAttemptSubmitted}, or \texttt{StreakUpdated} events are processed.
  \item \textbf{ResumeView.} A materialised row per learner recording the most recently accessed lesson, percentage through the video, and the last active timestamp, enabling efficient ``resume where you left off'' endpoints without scanning progress logs.
  \item \textbf{LessonAnalytics.} Aggregated view per lesson: total views, average completion rate, average quiz score, and discussion activity counts, computed from event streams rather than expensive OLAP queries on operational tables.
  \item \textbf{CourseSearchIndex.} A denormalised table containing course title, description, instructor name, category, average rating, and enrolment count, used by the search endpoint instead of joining across five write-optimised tables on every request.
\end{itemize}

These projections are stored in the same PostgreSQL instance but in a separate \texttt{read\_models} schema, reinforcing the conceptual separation between write-optimised domain tables and read-optimised projection tables. The outbox processor guarantees that projections eventually converge to consistency with the write model, typically within milliseconds of the source event being committed. Mehta et al.~\cite{mehta2023eda} have demonstrated that this eventual-consistency approach scales to educational platforms processing millions of learner interaction events per day, with projection latency dominated by the outbox polling interval rather than database throughput.

% =============================================================================
% Section 3.3: Background on AI in E-Learning
% =============================================================================

\section{Background on AI in E-Learning}
\label{sec:ai-elearning}

\subsection{OpenAI Whisper for Speech Recognition}

OpenAI Whisper~\cite{radford2023whisper} is a general-purpose speech recognition model trained on 680,000 hours of multilingual and multitask supervised data collected from the web. Unlike prior automatic speech recognition (ASR) systems that relied on meticulously curated, single-domain datasets, Whisper leverages weak supervision at an unprecedented scale, resulting in robust performance across diverse acoustic environments, speaker accents, and recording qualities. The model architecture is a sequence-to-sequence Transformer encoder-decoder: audio is resampled to 16 kHz, converted to an 80-channel log-mel spectrogram, and encoded by the Transformer encoder; the decoder then autoregressively produces the transcript text, optionally including timestamps, language identification, and segment boundaries.

Whisper's relevance to e-learning platforms is established by Panthel et al.~\cite{panthel2024whisper}, who evaluated Whisper transcriptions against human-produced subtitles across 120 university lecture recordings. They reported word error rates between 5.2\% and 8.1\% on English lectures with clear audio, rising to 12--15\% on recordings with background noise, accented speech, or specialised technical terminology. Critically, the study concluded that Whisper's output was ``sufficient for automated downstream tasks such as keyword extraction, content indexing, and AI quiz generation [...] with manual correction recommended only for public-facing subtitles where near-perfect accuracy is required.''

NexaLearn integrates Whisper through a dedicated \textbf{AI Pipeline} bounded context. When a video asset completes transcoding (Section~\ref{sec:video-streaming}), the pipeline:

\begin{enumerate}
  \item Downloads the MP4 rendition to a transient processing environment.
  \item Extracts the audio track and segments it into 30-second chunks aligned with the Whisper context window.
  \item Sends each chunk to the Whisper API or a self-hosted Whisper instance (configurable per deployment), requesting transcripts with word-level timestamps and segment boundaries.
  \item Assembles the per-chunk transcripts into a monolithic transcript document, preserving the temporal mapping between transcript segments and video offset.
  \item Stores the transcript as a JSON document associated with the \texttt{VideoAsset} entity, enabling downstream consumers to retrieve ``the transcript segment starting at second 145'' for fine-grained synchronisation.
\end{enumerate}

\subsection{LLM-Based Quiz Generation}

The second AI capability of NexaLearn is automated quiz generation from lecture transcripts using large language models. Kumar et al.~\cite{kumar2023llmquiz} demonstrated an end-to-end pipeline that inputs a lecture transcript and outputs a structured set of multiple-choice questions, true-or-false statements, and short-answer prompts, each with a correct answer, plausible distractors, and a justification. Their evaluation across 50 university lecture transcripts showed that GPT-4--generated questions were rated by domain experts as acceptable or better in 73\% of cases, with the primary failure mode being questions that tested recall of peripheral transcript details rather than conceptual understanding.

NexaLearn's AI pipeline generalises this approach. Upon receiving a transcript (either generated by Whisper or uploaded manually by the instructor), the pipeline constructs a prompt instructing the LLM to produce structured JSON adhering to a predefined schema:

\begin{lstlisting}[caption=AI-generated quiz JSON schema,label=lst:ai-quiz-schema]
{
  "questions": [
    {
      "stem": "What is the primary invariant enforced by ...?",
      "type": "multiple-choice",
      "options": [
        { "id": "A", "text": "...", "isCorrect": false },
        { "id": "B", "text": "...", "isCorrect": true  },
        { "id": "C", "text": "...", "isCorrect": false },
        { "id": "D", "text": "...", "isCorrect": false }
      ],
      "justification": "..."
    }
  ]
}
\end{lstlisting}

The generated quiz payload is not published directly to learners. Instead, NexaLearn applies a \textbf{human-in-the-loop review} workflow: the quiz is stored with a \texttt{READY\_FOR\_REVIEW} status, and the instructor is notified (via the Notifications context) to review, edit, approve, or reject each proposed question. Only upon instructor approval does the quiz status transition to \texttt{PUBLISHED}, making it available to enrolled learners. This review window ensures that the instructor retains pedagogical control while benefiting from automated question drafting---reducing the time required to create a quiz from a 60-minute lecture from approximately two hours of manual authoring to fifteen minutes of review and refinement.

\subsection{HMAC-Signed Callback Authentication}

The AI pipeline operates asynchronously: NexaLearn submits a transcription or quiz-generation request, the external AI service processes it (potentially taking seconds to minutes), and the service calls back a NexaLearn endpoint with the results. To authenticate these callbacks, NexaLearn implements an HMAC-signature scheme inspired by Stripe webhook signatures~\cite{stripe2024api}. Each callback request includes an \texttt{X-Pipeline-Signature} header with format \texttt{t=<unix-timestamp>,v1=<hex-digest>}. The recipient---the AI Pipeline controller---recomputes the HMAC-SHA256 digest of the request body using the shared secret, verifies that the timestamp is within a configurable tolerance window (default 300 seconds), and rejects requests whose signature does not match or whose timestamp is stale. This prevents replay attacks and ensures that only callbacks from the genuine AI pipeline are accepted.

The HMAC callback pattern, combined with the outbox event processing described in Section~\ref{sec:event-driven-outbox}, guarantees that AI-generated content is durably stored even if the application crashes during callback handling. The callback handler inserts the generated transcript or quiz into the database within a transaction that also creates an outbox event; the outbox processor subsequently notifies the instructor and updates related projections.

% =============================================================================
% Section 3.4: Background on Video Streaming Technology
% =============================================================================

\section{Background on Video Streaming Technology}
\label{sec:video-streaming}

\subsection{HTTP Live Streaming (HLS)}

HTTP Live Streaming (HLS), standardised as RFC 8216~\cite{pantos2022hls}, is the dominant adaptive-bitrate streaming protocol used by major video platforms including YouTube, Twitch, and Netflix. HLS works by segmenting a source video into short (typically 2--10 second) chunks encoded at multiple bitrates and resolutions. A \textbf{manifest file} (the \texttt{.m3u8} playlist) enumerates the available renditions and their segment URLs. The client---a web browser or mobile app---dynamically selects the highest bitrate segment that its current network conditions can sustain, switching to a lower bitrate when bandwidth drops and back to a higher bitrate when bandwidth recovers. This adaptive behaviour eliminates buffering interruptions while maximising playback quality.

The HLS specification defines a hierarchical playlist structure. A \textit{master playlist} lists all variant streams (e.g., 360p, 720p, 1080p) and auxiliary media such as subtitles and alternate audio tracks. Each variant stream has its own \textit{media playlist} containing the sequence of segment URIs and their durations. The protocol supports encryption via AES-128 or SAMPLE-AES, allowing content access control through key delivery authorisation.

\subsection{Mux Video Platform}

Mux~\cite{mux2024} is a video API platform that abstracts the complexity of video ingestion, transcoding, storage, and delivery. NexaLearn integrates with Mux through the \textbf{Media} bounded context, which manages the full lifecycle of video assets. Mux provides:

\begin{itemize}
  \item \textbf{Direct Uploads.} The client uploads a video file directly to Mux's storage via a \textit{direct upload endpoint}, obtaining a \texttt{upload\_id} and a playback ID. The upload is initiated by NexaLearn's backend, which requests a signed upload URL from Mux and returns it to the client; the client then PUTs the file directly to Mux, avoiding double-hop through the NexaLearn server.
  \item \textbf{Transcoding Pipeline.} Upon receiving a complete upload, Mux transcodes the source video into multiple HLS renditions (typically 360p, 480p, 720p, and 1080p) with configurable video codec (H.264), audio codec (AAC), and segment duration.
  \item \textbf{Webhook Notifications.} Mux sends webhooks to NexaLearn for lifecycle events: \texttt{video.upload.asset\_created}, \texttt{video.asset.ready}, \texttt{video.asset.errored}, \texttt{video.asset.deleted}. These webhooks drive the video asset state machine.
  \item \textbf{Playback IDs and Signed URLs.} Each transcoded asset receives a playback ID. For gated content, Mux supports \textit{signed playback URLs} using JSON Web Tokens (JWT)~\cite{jones2015jwt}. The NexaLearn backend generates a JWT containing the playback ID, an expiration time, and optional access restrictions (e.g., referrer whitelist, IP allowlist), signed with a Mux-specific signing key. The client embeds this signed JWT in the playback URL; Mux's delivery infrastructure verifies the signature and expiration before serving segments.
\end{itemize}

\subsection{Video Asset Lifecycle in NexaLearn}

The Media bounded context implements a state machine for each \texttt{VideoAsset} entity. The states are:

\begin{description}
  \item[\texttt{AWAITING\_UPLOAD}.] Initial state. A \texttt{Video} or \texttt{Lesson} entity has been created with an association to a new video, but the instructor has not yet uploaded the file. No Mux resource exists.
  \item[\texttt{UPLOADED}.] The instructor has initiated an upload and Mux has confirmed receipt of the source file. The state transitions from AWAITING\_UPLOAD to UPLOADED.
  \item[\texttt{TRANSCODING}.] Mux sends \texttt{video.upload.asset\_created}, indicating that the source file is being transcoded into HLS renditions. The NexaLearn backend creates a \texttt{GenerationJob} entity for the AI pipeline (Section~\ref{sec:ai-elearning}) in QUEUED state, awaiting the completed transcript.
  \item[\texttt{READY}.] Mux sends \texttt{video.asset.ready}, indicating that all HLS renditions are available. NexaLearn updates the \texttt{VideoAsset} with the playback ID and signed URL expiration policy. The AI pipeline receives a \texttt{VideoReady} event via the outbox, triggering the transcript-generation workflow.
  \item[\texttt{FAILED}.] Transcoding failed. NexaLearn logs the Mux error code and message, notifies the instructor, and makes the upload available for retry.
\end{description}

\subsection{MP4 Renditions for the AI Pipeline}

While HLS is the delivery format for learner playback, the Whisper transcription model (Section~\ref{sec:ai-elearning}) requires a contiguous audio or video file, not segmented HLS chunks. Mux provides a static MP4 rendition for each asset, accessible through a separate signed URL with an extended expiration window---sufficient for the AI pipeline to download the file, extract the audio track using FFmpeg, and feed it to Whisper. This dual-delivery strategy---HLS for learners, MP4 for automated processing---is transparent to the instructor, who uploads a single source file.

The interplay between video lifecycle and the AI pipeline exemplifies NexaLearn's event-driven design: a single external event (\texttt{video.asset.ready}) propagates through the outbox to trigger transcript generation, quiz creation, instructor notification, and search index updates, all without explicit orchestration code in the Media context. An asset may also transition directly to FAILED from UPLOADED or TRANSCODING if an error occurs at any stage.

% =============================================================================
% Section 3.5: Comparative Study of Previous Work
% =============================================================================

\section{Comparative Study of Previous Work}
\label{sec:comparative-study}

Several academic and industrial projects have explored architectural patterns for e-learning backends, event-driven reliability, and AI-assisted content generation. This section positions NexaLearn within the landscape of prior work by analysing four representative studies and identifying the gaps that NexaLearn's combined approach addresses.

\subsection{Dossena et al. (2022) --- DDD in E-Learning Platforms}

Dossena, Fugini, and Peroni~\cite{dossena2022ddd} present a case study of applying DDD strategic and tactical patterns to the design of a university LMS. Their work decomposes the e-learning domain into six bounded contexts---Course Management, Enrollment, Assessment, Gradebook, Communication, and Reporting---and documents the context map with upstream/downstream relationships. The authors report that the bounded context decomposition reduced inter-module coupling by 40\% compared to a previous monolithic version, and that the ubiquitous language improved communication between developers and faculty stakeholders during requirements gathering.

\textbf{Limitations.} The project was constrained to a single-semester pilot involving 15 courses and 300 students, limiting generalisability to production-scale deployments. The architecture relies on a relational database with no event-driven patterns, meaning that cross-context synchronisation is performed through direct remote procedure calls (RPC) between microservices, exposing the system to the dual-write problem described in Section~\ref{sec:event-driven-outbox}. The assessment context supports only manually authored quizzes, with no AI-assisted question generation.

\subsection{Yang et al. (2023) --- Microservices for E-Learning}

Yang, Chen, and Zhang~\cite{yang2023micro} conducted a systematic literature review of 47 papers published between 2018 and 2023 on microservice decomposition, inter-service communication, and data management in e-learning systems. Their taxonomy identifies four decomposition strategies: (1) by business capability (e.g., separate services for users, courses, payments), (2) by workflow step (e.g., video ingestion as a separate service from video delivery), (3) by subdomain (consistent with DDD bounded contexts), and (4) by volatility (extracting frequently changing features into independently deployable services). The survey found that 62\% of analysed systems used synchronous HTTP/REST for inter-service communication, 21\% used asynchronous messaging (primarily RabbitMQ or Apache Kafka), and only 17\% employed transactional outbox patterns for data consistency.

\textbf{Limitations.} The review identifies a significant adoption gap: while the outbox pattern is recommended in the microservices literature ~\cite{richardson2018outbox} as the standard solution for reliable event publication, fewer than one in five surveyed e-learning systems implement it. Most rely on distributed transactions (XA protocol) or best-effort eventual consistency without idempotency guarantees. The review does not address AI integration patterns, as none of the surveyed papers incorporated Whisper, LLM-based quiz generation, or similar AI services as first-class components of the architecture.

\subsection{Kumar et al. (2023) --- LLM Quiz Generation}

Kumar, Barrios, and Haridas~\cite{kumar2023llmquiz} present an end-to-end pipeline for generating formative assessment questions from lecture transcripts using GPT-4. Their system, evaluated on 50 university lectures spanning computer science, biology, and economics, achieved a 73\% expert-acceptability rate for generated multiple-choice questions. The paper contributes a prompt engineering methodology that instructs the LLM to produce questions targeting different cognitive levels (Bloom's taxonomy: remember, understand, apply, analyse).

\textbf{Limitations.} The pipeline is a standalone Python service invoked manually by instructors; it does not integrate with an LMS backend for automated lifecycle management. There is no support for human-in-the-loop review workflows, meaning that generated questions are either used directly (risk of uncaught errors) or manually copied into a separate quiz-authoring tool. The paper does not address authentication or secure callback patterns for async AI processing, and the pipeline assumes the transcript is already available---it does not include ASR integration.

\subsection{Mehta et al. (2023) --- Event-Driven Educational Analytics}

Mehta, Soni, and Thakkar~\cite{mehta2023eda} propose an event-driven architecture for processing real-time learner interaction events in large-scale educational platforms. Their system defines a canonical event schema---\texttt{LearnerInteractionEvent} with fields for learner ID, session ID, resource ID, event type (e.g., \texttt{PAGE\_VIEW}, \texttt{VIDEO\_PLAY}, \texttt{QUIZ\_SUBMIT}), timestamp, and arbitrary key-value metadata. Events are ingested via Apache Kafka and processed by Kafka Streams topologies that compute per-learner and per-course aggregates with sub-second latency. The authors report processing 2.8 million events per day across 15,000 learners with a p99 processing latency of 850 ms.

\textbf{Limitations.} The architecture depends on a Kafka cluster, which represents significant operational overhead for self-hosted deployments by small teams or institutions with limited DevOps resources. The event schema is generic and does not enforce domain-specific invariants; events are not derived from domain aggregates and therefore risk becoming disconnected from the write model's correctness guarantees. The outbox pattern is not used, as events are published directly to Kafka from application code, exposing the system to the dual-write problem.

\subsection{Synthesis and Gap Analysis}

Table~\ref{tab:lit-comparison} summarises the coverage of key architectural dimensions across the four studies and NexaLearn using $\checkmark$ (full coverage), $\sim$ (partial/reviewed), and $\times$ (not addressed). The comparison reveals that no single prior work addresses the full combination of DDD decomposition, reliable event processing via the outbox pattern, AI-assisted transcription and quiz generation with HMAC-secured callbacks, and HLS-based video streaming with lifecycle management. NexaLearn's contribution is the \textit{integration} of these patterns into a cohesive, open-source backend architecture that is deployable as a modular monolith without requiring a separate message broker infrastructure.

\begin{table}[H]
  \centering
  \caption{Comparative Summary of Prior Work and NexaLearn Coverage}
  \label{tab:lit-comparison}
  \small
  \begin{tabularx}{\textwidth}{@{}X c c c c c@{}}
    \toprule
    \textbf{Dimension} &
    \textbf{Dossena} &
    \textbf{Yang} &
    \textbf{Kumar} &
    \textbf{Mehta} &
    \textbf{NexaLearn} \\
    \midrule
    DDD bounded contexts     & $\checkmark$ & $\sim$ & $\times$ & $\times$ & $\checkmark$ \\
    \addlinespace
    Clean Architecture layers & $\sim$ & $\times$ & $\times$ & $\times$ & $\checkmark$ \\
    \addlinespace
    Transactional outbox     & $\times$ & $\sim$ & $\times$ & $\times$ & $\checkmark$ \\
    \addlinespace
    Idempotent event processing & $\times$ & $\times$ & $\times$ & $\sim$ & $\checkmark$ \\
    \addlinespace
    ASR (Whisper) integration & $\times$ & $\times$ & $\times$ & $\times$ & $\checkmark$ \\
    \addlinespace
    LLM quiz generation      & $\times$ & $\times$ & $\checkmark$ & $\times$ & $\checkmark$ \\
    \addlinespace
    Human-in-the-loop review & $\times$ & $\times$ & $\times$ & $\times$ & $\checkmark$ \\
    \addlinespace
    HLS video streaming      & $\times$ & $\sim$ & $\times$ & $\times$ & $\checkmark$ \\
    \addlinespace
    HMAC callback security   & $\times$ & $\times$ & $\times$ & $\times$ & $\checkmark$ \\
    \addlinespace
    No broker dependency     & $\times$ & $\sim$ & $\times$ & $\times$ & $\checkmark$ \\
    \midrule
    \textbf{Rationale} &
    \parbox[t]{2.2cm}{\footnotesize DDD-focused; no AI or eventing} &
    \parbox[t]{2.2cm}{\footnotesize Survey only; no implementation} &
    \parbox[t]{2.2cm}{\footnotesize Standalone LLM; no backend integration} &
    \parbox[t]{2.2cm}{\footnotesize Kafka-based; no DDD or AI} &
    \parbox[t]{2.2cm}{\footnotesize \textbf{All dimensions integrated with modular monolith}} \\
    \bottomrule
  \end{tabularx}
\end{table}

NexaLearn is distinguished by its comprehensive coverage across all ten dimensions. The three dimensions not covered by any prior work---Whisper integration, human-in-the-loop quiz review, and HMAC callback security---represent the principal novel contributions of this project's architecture. The absence of a mandatory message broker (Kafka or RabbitMQ) is a deliberate design decision, justified in Section~\ref{sec:implemented-approach}, that reduces operational complexity for the primary target audience of edtech startups and university IT departments.

The table further reveals that existing work either addresses architecture (Dossena, Yang, Mehta) or AI content (Kumar), but never both. NexaLearn's architecture is explicitly designed to bridge this gap, making it---to the best of our knowledge---the first open-source NestJS backend to combine DDD-governed bounded contexts with a full AI pipeline, transactional outbox eventing, and Mux-based video lifecycle.

% =============================================================================
% Section 3.6: Implemented Approach and Justification
% =============================================================================

\section{Implemented Approach and Justification}
\label{sec:implemented-approach}

This section synthesises the literature surveyed in Sections~\ref{sec:ddd-clean-arch} through~\ref{sec:comparative-study} into concrete technology and pattern decisions for NexaLearn. Each decision is justified against alternatives on technical merit, alignment with the problem statement (Chapter~1), and suitability for the primary target segments identified in Section~\ref{sec:targeted-customers}.

\subsection{Framework Selection: NestJS over Spring Boot and Django}

NexaLearn is implemented in \textbf{NestJS}\cite{nestjs2024}, a progressive Node.js framework that leverages TypeScript decorators for dependency injection, controller routing, middleware piping, and module organisation. Three considerations motivate this choice:

\begin{enumerate}
  \item \textbf{TypeScript ecosystem alignment.} The primary target segment---edtech startups and independent developers---operates predominantly in the TypeScript/JavaScript ecosystem. Choosing NestJS allows these developers to read, fork, and contribute to NexaLearn without learning a second programming language. Garcia et al.~\cite{garcia2022typescript} found that TypeScript adoption reduces production defect density by 15\% compared to equivalent JavaScript codebases, primarily through compile-time detection of type mismatches, null reference errors, and incorrect API usage.
  \item \textbf{Decorator-based DDD mapping.} NestJS's decorator syntax maps naturally to DDD architectural layers. A controller is annotated with \texttt{@Controller()}, a provider with \texttt{@Injectable()}, and a module with \texttt{@Module()}. The framework provides no prescribed entity or aggregate base class---NexaLearn's domain entities are plain TypeScript classes---but the dependency injection container cleanly separates application-layer use cases from infrastructure adapters.
  \item \textbf{Comparative framework analysis.} Zhang et al.~\cite{zhang2024nestjs} conducted a controlled experiment comparing NestJS, Spring Boot, and Django on throughput, memory consumption, and developer productivity for a representative CRUD+event API workload. NestJS achieved comparable throughput to Spring Boot (within 12\%) with 40\% lower memory consumption in containerised deployments, and developers with TypeScript experience completed identical implementation tasks 1.7\texttimes{} faster in NestJS than developers with Java experience in Spring Boot.
\end{enumerate}

\subsection{ORM Selection: Prisma over TypeORM}

Prisma~\cite{prisma2024} was chosen as the object-relational mapping layer over TypeORM and MikroORM based on three criteria:

\begin{enumerate}
  \item \textbf{Type safety.} Prisma generates TypeScript types from the database schema declaration (\texttt{prisma/schema.prisma}). Every query returns a fully typed result set, eliminating the ``field does not exist on type'' runtime errors common with TypeORM's class-based entity approach. The generated types serve as a single source of truth for schema evolution.
  \item \textbf{Migration management.} Prisma Migrate produces deterministic PostgreSQL migration files with rollback support, integrated into the development workflow via \texttt{prisma migrate dev} and deployment pipeline via \texttt{prisma migrate deploy}. TypeORM migrations, by contrast, require careful ordering and manual conflict resolution in team environments.
  \item \textbf{Raw query support.} While Prisma's generated client covers 95\% of NexaLearn's data access patterns, the CQRS-lite read models (Section~\ref{subsec:revenue}) and analytics aggregations require window functions, lateral joins, and materialised view refreshes that benefit from raw SQL. Prisma exposes \texttt{prisma.\$queryRawUnsafe()} with parameterised inputs, allowing these queries without escaping the typed ORM layer.
\end{enumerate}

\subsection{Transactional Outbox over Direct Message Queue Publication}

NexaLearn implements the transactional outbox pattern (Section~\ref{sec:event-driven-outbox}) using PostgreSQL as the event store, with no additional message broker such as RabbitMQ or Apache Kafka. This decision is motivated by the analysis of Mehta et al.~\cite{mehta2023eda} and the gap identified in Table~\ref{tab:lit-comparison}:

\begin{enumerate}
  \item \textbf{Deployment simplicity.} A PostgreSQL-only dependency means that NexaLearn deploys as a single Docker Compose stack (application containers, PostgreSQL, Redis for caching) rather than requiring a Kafka or RabbitMQ cluster. This is a first-class requirement for the primary customer segment: university IT departments and bootstrapped startups that may not have dedicated DevOps engineers to operate a distributed messaging infrastructure.
  \item \textbf{Transactional guarantees.} The outbox pattern ensures that events are persisted atomically with the domain mutation, solving the dual-write problem without distributed transactions. A message broker introduces a separate transactional boundary: the application would need to write to the database and then publish to the broker, with the risk that one succeeds while the other fails.
  \item \textbf{Acceptable latency.} NexaLearn's polling interval defaults to 500 ms with a batch size of 50 events. For the workload profile of an e-learning platform serving primarily synchronous (learner-facing) and asynchronous (instructor notification, search index update) consumers, this latency is well within acceptable bounds. Events such as \texttt{LessonCompleted} that trigger real-time progress UI updates are additionally written to Redis for immediate consumption, combining the outbox's durability with Redis's pub/sub latency.
\end{enumerate}

The trade-off is increased database load from polling and the lack of native fan-out, message broker, delayed-queue, or dead-letter features. For NexaLearn's current workload targets (tens of thousands of enrolled learners, hundreds of events per second at peak), PostgreSQL with appropriate indexing (\texttt{idx\_outbox\_unprocessed} on \texttt{(processed\_at NULLS FIRST, created\_at ASC)}) comfortably handles this volume. Should event throughput grow beyond this bound, the outbox publisher is designed to be replaced with a Kafka producer without changing any domain or application layer code---the outbox processor's only external contract is a \texttt{EventPublisher} interface that currently calls a local NestJS event bus but could be reimplemented to publish to Kafka.

\subsection{Modular Monolith over Microservices}

NexaLearn is deployed as a \textbf{modular monolith}: a single NestJS application process with multiple modules, rather than a distributed system of independently deployed microservices. The modular monolith is the deployment model recommended by Yang et al.~\cite{yang2023micro} for projects and teams at NexaLearn's scale, offering:

\begin{enumerate}
  \item \textbf{Reduced operational overhead.} A single application instance is simpler to deploy, monitor, and debug than sixteen microservices with their own deployment pipelines, health checks, and log aggregation. For a graduation project maintained by a small core team, minimising operational surface area is a practical necessity.
  \item \textbf{No network overhead.} Inter-context communication within the same process uses in-memory method calls rather than HTTP/gRPC round-trips, reducing latency by several orders of magnitude for high-frequency interactions such as ``validate enrollment before returning lesson content.''
  \item \textbf{Explicit bounded contexts.} The modular monolith does not relax DDD boundaries. Modules cannot import repository classes from other modules; cross-context communication is mediated through domain events (published via the in-process event bus and the outbox) or through dedicated query services exposed as NestJS providers registered in the importing module. This discipline ensures that extracting any context into an independent microservice---should scaling requirements eventually demand it---requires changing only the module's registration and the event publishing mechanism, not the domain logic.
\end{enumerate}

\subsection{Summary of Architectural Decisions}

Table~\ref{tab:arch-decisions} summarises the key architectural decisions, the alternatives considered, and the primary justification for each choice.

\begin{table}[H]
  \centering
  \caption{Architectural Decisions and Justification}
  \label{tab:arch-decisions}
  \small
  \begin{tabular}{@{}p{3.0cm} p{2.5cm} p{2.5cm} p{4.5cm}@{}}
    \toprule
    \textbf{Decision} &
    \textbf{Chosen} &
    \textbf{Alternatives} &
    \textbf{Primary Justification} \\
    \midrule
    Framework &
    NestJS &
    Spring Boot, Django &
    TypeScript ecosystem alignment with target segment; decorator-based DI; lower memory in containers \\
    \addlinespace
    ORM &
    Prisma &
    TypeORM, MikroORM &
    Generated type safety; deterministic migrations; raw query escape hatch \\
    \addlinespace
    Event transport &
    Outbox (PostgreSQL) &
    Kafka, RabbitMQ &
    No broker dependency for simplified deployment; atomic dual-write guarantees \\
    \addlinespace
    Deployment model &
    Modular monolith &
    Microservices &
    Reduced ops overhead; no network latency; extractable to microservices later \\
    \addlinespace
    AI pipeline invocation &
    HMAC-signed async callback &
    Polling, synchronous API &
    Stripe-style security; decoupled processing; built-in replay prevention \\
    \addlinespace
    Video delivery &
    Mux (HLS + signed URLs) &
    Cloudflare Stream, self-hosted FFmpeg &
    Managed transcoding; JWT-signed playback; webhook lifecycle events \\
    \bottomrule
  \end{tabular}
\end{table}

These decisions, grounded in the literature surveyed throughout this chapter, collectively deliver an open-source e-learning backend that balances architectural rigour (DDD, Clean Architecture, outbox), practical deployability (modular monolith, PostgreSQL-only event store), and modern AI integration (Whisper, LLM quiz generation, HMAC security)---the three pillars that define NexaLearn's contribution relative to the prior work examined in Section~\ref{sec:comparative-study}.

